\documentclass{beamer}

\mode<presentation>
\title{SoC1\&2: Sign-off Timing with Quantus and Tempus}
\author[AS]{Andreas Siggelkow \newline \texttt{andreas.siggelkow@rwu.de}}
\date[2026]{\today}
\subtitle{RWU}
\institute{Elektrotechnik und Informatik}

%%%%%%%%%%%%%%%%%%%%%%%%%%%%%%%%
% Logo
%%%%%%%%%%%%%%%%%%%%%%%%%%%%%%%%
\pgfdeclareimage[height=1cm]{hslogo}{./hs}
\logo{\pgfuseimage{hslogo}}
%%%%%%%%%%%%%%%%%%%%%%%%%%%%%%%%
% Macros
%%%%%%%%%%%%%%%%%%%%%%%%%%%%%%%%
\usepackage[T1]{fontenc}
\usepackage{lmodern}
\usepackage{textcomp}
\definecolor{links}{HTML}{2A1B81}
\hypersetup{colorlinks,linkcolor=,urlcolor=links}
\usepackage{graphicx}
\usepackage{ragged2e, tabularx}
\newcolumntype{Y}{>{\RaggedRight\arraybackslash}X}
\usetheme{Singapore}
\usecolortheme{seahorse}
\usefonttheme[onlylarge]{structuresmallcapsserif}
\usefonttheme[onlysmall]{structurebold}
\setbeamertemplate{footline}[frame number]
\usepackage{listings}
\newcommand{\astitle}{nix}
\newcommand{\assubtitle}{nix}
\newcommand{\asfolie}[3]{\begin{frame}\frametitle{#1}\begin{block}{#2}{#3}\end{block}\end{frame}}

% Commands and tool output are set with listings: dashes, $, _, \ and quotes
% appear exactly as they have to be typed.
\definecolor{lstbg}{HTML}{F2F2F7}
\definecolor{lstcomment}{HTML}{5A5A8A}
\lstset{
  basicstyle=\ttfamily\scriptsize,
  columns=fullflexible,
  keepspaces=true,
  upquote=true,
  breaklines=true,
  breakatwhitespace=false,
  backgroundcolor=\color{lstbg},
  frame=none,
  xleftmargin=2pt,
  xrightmargin=2pt,
  aboveskip=4pt,
  belowskip=4pt,
  showstringspaces=false
}
\lstdefinestyle{sh}{
  morecomment=[l]{\#},
  commentstyle=\color{lstcomment}\itshape
}
\lstdefinestyle{out}{basicstyle=\ttfamily\tiny}
\newcommand{\cmd}[1]{\texttt{#1}}



\begin{document}
%%%%%%%%%%%%%%%%%%%%%%%%%%%%%%%%
% First pages
%%%%%%%%%%%%%%%%%%%%%%%%%%%%%%%%
\setcounter{tocdepth}{1}

\begin{frame}
  \titlepage
\end{frame}

\begin{frame}
  \frametitle{Outline}
  \tableofcontents
\end{frame}

%%%%%%%%%%%%%%%%%%%%%%%%%%%%%%%%
% Overview
%%%%%%%%%%%%%%%%%%%%%%%%%%%%%%%%
\section{Overview}
\renewcommand{\astitle}{Overview}

\renewcommand{\assubtitle}{Where in the Flow?}
\begin{frame}
\frametitle{\astitle}
\begin{block}{\assubtitle}
  \footnotesize
  \begin{tabularx}{\linewidth}{@{}llY@{}}
    \textbf{step} & \textbf{tool / script} & \textbf{result}\\\hline
    RTL simulation      & xrun, \cmd{scripts/xcelium/}   & functionally correct RTL\\
    synthesis + DFT     & Genus, \cmd{scripts/genus/}    & netlist, scan chain\\
    gate-level sim      & xrun, \cmd{scripts/gatesim/}   & netlist ok\\
    place \& route      & Innovus, \cmd{scripts/innovus/} & layout, \cmd{\_routed.v/.def}\\
    \textbf{RC extraction} & \textbf{Quantus}, \cmd{scripts/quantus/} & SPEF per RC corner\\
    \textbf{sign-off STA}  & \textbf{Tempus}, \cmd{scripts/sta/}      & slack per view, SDF\\
    post-route GLS      & xrun, \cmd{gatesim.sh -p -s}   & netlist with real delays\\
    physical verification & Pegasus/Calibre             & DRC, LVS, ERC\\
  \end{tabularx}
  \vspace{4pt}
  \begin{itemize}
    \item This deck: the two bold steps, always in this order
  \end{itemize}
\end{block}
\end{frame}

\renewcommand{\assubtitle}{Why Extraction First?}
\begin{frame}
\frametitle{\astitle}
\begin{block}{\assubtitle}
  \begin{itemize}
    \item A path delay is cell delay + wire delay
    \item Tempus computes the cell delays from the Liberty tables, but it
      does \emph{not} know the wires: their resistance and capacitance
      come from the layout
    \item Quantus reads the routed layout (DEF + LEF) and the process
      description (\cmd{qrcTechFile}) and writes the RC network of every
      net into a \emph{SPEF} file
    \item Without SPEF Tempus can only estimate the wires -- that is what
      Genus does before the layout, not a sign-off
    \item The timing reports inside Innovus use a fast internal
      extraction; good for optimization, not accurate enough for sign-off
  \end{itemize}
\end{block}
\end{frame}

\renewcommand{\assubtitle}{Data Flow}
\begin{frame}
\frametitle{\astitle}
\begin{block}{\assubtitle}
  \footnotesize
  \begin{tabularx}{\linewidth}{@{}lY@{}}
    \textbf{file} & \textbf{from $\rightarrow$ to}\\\hline
    \cmd{cc\_chip\_top\_routed.def} & Innovus $\rightarrow$ Quantus (layout)\\
    \cmd{cc\_chip\_top\_routed.v}   & Innovus $\rightarrow$ Tempus, GLS (netlist with clock tree)\\
    LEF, \cmd{qrcTechFile} per corner & PDK $\rightarrow$ Quantus\\
    \cmd{cc\_chip\_top\_rc\_\{max,typ,min\}.spef} & Quantus $\rightarrow$ Tempus\\
    Liberty per corner, SDC         & PDK, \cmd{scripts/} $\rightarrow$ Tempus\\
    \cmd{cc\_chip\_top\_sta\_\{func,scan\}.sdf} & Tempus $\rightarrow$ GLS\\
  \end{tabularx}
  \vspace{4pt}
  \begin{itemize}
    \item Build directories: \cmd{build\_innovus/}, \cmd{build\_quantus/},
      \cmd{build\_sta/} -- all in \cmd{.gitignore} (X-FAB NDA)
  \end{itemize}
\end{block}
\end{frame}

%%%%%%%%%%%%%%%%%%%%%%%%%%%%%%%%
% Setup
%%%%%%%%%%%%%%%%%%%%%%%%%%%%%%%%
\section{Setup}
\renewcommand{\astitle}{Setup}

\renewcommand{\assubtitle}{Prerequisites}
\begin{frame}[fragile]
\frametitle{\astitle}
\begin{block}{\assubtitle}
  \begin{itemize}
    \item Cadence environment with \cmd{innovus}, \cmd{quantus},
      \cmd{tempus} in the \cmd{PATH}
    \item Place \& route done:
\begin{lstlisting}[style=sh]
cd RWU_RV/RV_NoPipeline
./scripts/innovus/innovus_pnr.sh -b     # -> build_innovus/
\end{lstlisting}
    \item \cmd{\$XFAB\_PATHS} (X-FAB NDA, never in git) already holds what
      Genus needs; used here:
\begin{lstlisting}[style=sh]
set LEF_FILES [list <tech.lef> <std_cells.lef> ...]
set QRC_MAX <QRC-Max/qrcTechFile>     ;# also _TYP, _MIN
set LIBS_SLOW [list ...]              ;# also _TYP, _FAST
set QRC_LAYERMAP [list]   ;# new, optional: LEF/QRC layer pairs
\end{lstlisting}
  \end{itemize}
\end{block}
\end{frame}

\renewcommand{\assubtitle}{Hand-off from Innovus}
\begin{frame}
\frametitle{\astitle}
\begin{block}{\assubtitle}
  \begin{itemize}
    \item \cmd{02\_rv64\_floorplan\_pnr.tcl} ends with
      \cmd{02b\_rv64\_export.tcl}:
    \begin{itemize}
      \item \cmd{write\_netlist cc\_chip\_top\_routed.v}: logic incl. the
        CTS buffers, without fillers and corner cells
      \item \cmd{write\_def -floorplan -netlist -routing
        cc\_chip\_top\_routed.def}
    \end{itemize}
    \item P\&R ran before this step existed? No new P\&R run needed:
      \cmd{quantus\_rcx.sh} finds \cmd{cc\_chip\_top\_routed.db} and runs
      \cmd{02b\_rv64\_export.tcl} on it
    \item The netlist must be the routed one: the Genus netlist has no
      clock tree, its clock skew would be wrong
  \end{itemize}
\end{block}
\end{frame}

%%%%%%%%%%%%%%%%%%%%%%%%%%%%%%%%
% Quantus
%%%%%%%%%%%%%%%%%%%%%%%%%%%%%%%%
\section{RC Extraction}
\renewcommand{\astitle}{RC Extraction}

\renewcommand{\assubtitle}{RC Corners}
\begin{frame}
\frametitle{\astitle}
\begin{block}{\assubtitle}
  \footnotesize
  \begin{tabularx}{\linewidth}{@{}lllY@{}}
    \textbf{RC corner} & \textbf{QRC file} & \textbf{temp.} & \textbf{used with}\\\hline
    \cmd{rc\_max} & \cmd{QRC\_MAX} & 125\,$^\circ$C & slow cells: setup (\cmd{dc\_slow})\\
    \cmd{rc\_typ} & \cmd{QRC\_TYP} & 25\,$^\circ$C  & typical, reported only\\
    \cmd{rc\_min} & \cmd{QRC\_MIN} & $-40$\,$^\circ$C & fast cells: hold (\cmd{dc\_fast})\\
  \end{tabularx}
  \vspace{4pt}
  \begin{itemize}
    \item Same names as \cmd{create\_rc\_corner} in
      \cmd{xh018\_mmmc.tcl}: Tempus maps each SPEF to its corner by name
    \item Metal thickness, width and dielectric vary with the process: max
      = largest RC, min = smallest
    \item Extraction type \cmd{rc\_coupled}: the coupling capacitances
      between neighbouring wires are kept separately -- Tempus needs them
      for crosstalk (SI)
  \end{itemize}
\end{block}
\end{frame}

\renewcommand{\assubtitle}{Standalone Quantus: Generated Files}
\begin{frame}
\frametitle{\astitle}
\begin{block}{\assubtitle}
  \begin{itemize}
    \item Quantus reads a command file (CCL), not Tcl;
      \cmd{03a\_rv64\_quantus.tcl} runs in \cmd{tclsh} and writes into
      \cmd{build\_quantus/}:
    \begin{itemize}
      \item \cmd{qrc/rc\_max/qrcTechFile} etc.: links to the X-FAB files
      \item \cmd{lib.defs}: one line \cmd{DEFINE xh018 <dir>/qrc}, the
        corners are the subdirectories
      \item \cmd{quantus.ccl}: the run
    \end{itemize}
    \item The links rename the X-FAB corners to \cmd{rc\_max}\ldots: the
      SPEF files come out as \cmd{cc\_chip\_top\_rc\_max.spef}, no X-FAB
      name in any script
    \item The generated files hold X-FAB paths: build directory only
  \end{itemize}
\end{block}
\end{frame}

\renewcommand{\assubtitle}{Standalone Quantus: \texttt{quantus.ccl}}
\begin{frame}[fragile]
\frametitle{\astitle}
\begin{block}{\assubtitle}
\begin{lstlisting}[style=sh]
input_db -type def \
    -design_file "<...>/cc_chip_top_routed.def" \
    -lef_file_list "<tech.lef>" "<std_cells.lef>" ...
process_technology \
    -technology_library_file "<...>/lib.defs" \
    -technology_name "xh018" \
    -technology_corner "rc_max" "rc_typ" "rc_min" \
    -temperature 125 25 -40
extract \
    -selection "all" \
    -type "rc_coupled"
output_setup \
    -file_name "cc_chip_top" \
    -compressed false
output_db -type spef
\end{lstlisting}
  \scriptsize
  \begin{itemize}
    \item With \cmd{QRC\_LAYERMAP} set, an \cmd{extraction\_setup
      -technology\_layer\_map} block is added
  \end{itemize}
\end{block}
\end{frame}

\renewcommand{\assubtitle}{Alternative: Quantus Engine inside Innovus}
\begin{frame}[fragile]
\frametitle{\astitle}
\begin{block}{\assubtitle}
  \begin{itemize}
    \item \cmd{03a\_rv64\_rcx\_innovus.tcl}, option \cmd{-i}:
\begin{lstlisting}[style=sh]
read_db cc_chip_top_routed.db
set_db extract_rc_engine       post_route
set_db extract_rc_effort_level signoff   ;# = Quantus
set_db extract_rc_coupled      true
extract_rc
foreach rc {rc_max rc_typ rc_min} {
  write_parasitics -rc_corner $rc -spef_file cc_chip_top_${rc}.spef
}
\end{lstlisting}
    \item Takes the QRC files from the MMMC of the Innovus DB: no CCL, no
      \cmd{lib.defs}, no layer map
    \item Same SPEF names: Tempus does not care which way was used
    \item Use it if the standalone run fails, or to compare both
  \end{itemize}
\end{block}
\end{frame}

\renewcommand{\assubtitle}{Running the Extraction}
\begin{frame}[fragile]
\frametitle{\astitle}
\begin{block}{\assubtitle}
\begin{lstlisting}[style=sh]
./scripts/quantus/quantus_rcx.sh        # standalone Quantus
./scripts/quantus/quantus_rcx.sh -i     # engine inside Innovus
./scripts/quantus/quantus_rcx.sh -c     # clean build_quantus/ first
\end{lstlisting}
\begin{lstlisting}[style=out]
=== [quantus] 1/3 hand-off from .../build_innovus
=== [quantus] 2/3 standalone Quantus in .../build_quantus
=== [quantus] 3/3 result
  cc_chip_top_rc_max.spef     NNNNN nets
  cc_chip_top_rc_typ.spef     NNNNN nets
  cc_chip_top_rc_min.spef     NNNNN nets
PASS  RC extraction of cc_chip_top (quantus), next: ./scripts/sta/sta.sh -b
\end{lstlisting}
  \begin{itemize}
    \item Check: the three net counts are equal, no \cmd{ERROR} in
      \cmd{quantus.log} (\cmd{rcx.log} with \cmd{-i})
  \end{itemize}
\end{block}
\end{frame}

%%%%%%%%%%%%%%%%%%%%%%%%%%%%%%%%
% Tempus
%%%%%%%%%%%%%%%%%%%%%%%%%%%%%%%%
\section{Static Timing Analysis}
\renewcommand{\astitle}{Static Timing Analysis}

\renewcommand{\assubtitle}{Constraint Modes}
\begin{frame}
\frametitle{\astitle}
\begin{block}{\assubtitle}
  \footnotesize
  \begin{tabularx}{\linewidth}{@{}llY@{}}
    \textbf{mode} & \textbf{SDC} & \textbf{situation}\\\hline
    \cmd{func}      & \cmd{as\_top\_mem\_func.sdc} & CPU runs on \cmd{clk\_core} (1.25\,MHz)\\
    \cmd{jtag\_load} & \cmd{as\_top\_mem\_jtag.sdc} & JTAG loads the I-Mem, macro on \cmd{tck\_i}\\
    \cmd{scan}      & \cmd{as\_top\_mem\_scan.sdc} & SCANTEST: core on \cmd{tck\_i} (25\,MHz)\\
  \end{tabularx}
  \vspace{4pt}
  \begin{itemize}
    \item All three on top of \cmd{as\_top\_mem.sdc} (clocks, I/O, false
      paths); each mode file fixes \cmd{sc01\_mode\_o} and
      \cmd{im\_mode\_o} with \cmd{set\_case\_analysis}
    \item \cmd{scan} is new and only in the STA
      (\cmd{xh018\_mmmc\_sta.tcl}): Genus must not upsize the core logic
      for a 40\,ns capture that is not needed in operation
  \end{itemize}
\end{block}
\end{frame}

\renewcommand{\assubtitle}{Analysis Views}
\begin{frame}
\frametitle{\astitle}
\begin{block}{\assubtitle}
  \footnotesize
  \begin{tabularx}{\linewidth}{@{}lllY@{}}
    \textbf{view} & \textbf{mode} & \textbf{corner} & \textbf{check}\\\hline
    \cmd{av\_slow}      & func      & slow + rc\_max & setup\\
    \cmd{av\_typ}       & func      & typ + rc\_typ  & setup (info)\\
    \cmd{av\_jtag\_slow} & jtag\_load & slow + rc\_max & setup\\
    \cmd{av\_scan\_slow} & scan      & slow + rc\_max & setup\\
    \cmd{av\_fast}      & func      & fast + rc\_min & hold\\
    \cmd{av\_jtag\_fast} & jtag\_load & fast + rc\_min & hold\\
    \cmd{av\_scan\_fast} & scan      & fast + rc\_min & hold\\
  \end{tabularx}
  \vspace{4pt}
  \begin{itemize}
    \item Setup = data too late: slow cells, high RC, 125\,$^\circ$C
    \item Hold = data too early: fast cells, low RC, $-40$\,$^\circ$C
    \item Tempus analyses all views in one run
  \end{itemize}
\end{block}
\end{frame}

\renewcommand{\assubtitle}{Why a Scan Mode?}
\begin{frame}
\frametitle{\astitle}
\begin{block}{\assubtitle}
  \begin{itemize}
    \item In SCANTEST all scan flip-flops run on \cmd{tck\_i}, through
      \cmd{scan\_ckgate} and \cmd{core\_ckmux}
    \item \textbf{Shift}: flip-flop $\rightarrow$ flip-flop with no logic
      in between. The data path is short, so clock skew alone can break
      \emph{hold} -- the classic scan problem
    \item \textbf{Capture}: the functional logic, at 25\,MHz instead of
      1.25\,MHz
    \item Rules for \cmd{av\_scan\_*}:
    \begin{itemize}
      \item hold violations: must be 0 (the chain would not shift)
      \item setup violations: only limit the TCK of SCANTEST; the tester
        may capture with a slower clock
    \end{itemize}
  \end{itemize}
\end{block}
\end{frame}

\renewcommand{\assubtitle}{\texttt{03b\_rv64\_sta.tcl}}
\begin{frame}[fragile]
\frametitle{\astitle}
\begin{block}{\assubtitle}
\begin{lstlisting}[style=sh]
read_mmmc    $SCRIPT_DIR/xh018_mmmc_sta.tcl
read_netlist $NETLIST -top cc_chip_top
init_design
foreach rc {rc_max rc_typ rc_min} {
  read_spef -rc_corner $rc $SPEF_DIR/cc_chip_top_${rc}.spef
}
set_db timing_analysis_type ocv      ;# on-chip variation
set_db timing_analysis_cppr both     ;# common path pessimism removal
set_db delaycal_enable_si   true     ;# crosstalk (--no-si: off)
check_timing -verbose    > sta_check_timing.rpt
report_timing -late  -view av_slow -max_paths 20 > sta_setup_av_slow.rpt
report_timing -early -view av_fast -max_paths 20 > sta_hold_av_fast.rpt
report_constraint -all_violators > sta_violators.rpt
write_sdf -min_view av_fast -max_view av_slow cc_chip_top_sta_func.sdf
\end{lstlisting}
  \scriptsize
  \begin{itemize}
    \item \cmd{report\_timing} runs for every setup and hold view
  \end{itemize}
\end{block}
\end{frame}

\renewcommand{\assubtitle}{OCV, CPPR, SI}
\begin{frame}
\frametitle{\astitle}
\begin{block}{\assubtitle}
  \begin{itemize}
    \item \textbf{OCV}: launch and capture clock path may be on
      different parts of the die; Tempus takes the late delay for one and
      the early delay for the other
    \item \textbf{CPPR}: up to the branch point both clock paths run
      through the \emph{same} buffers -- that part cannot be early and
      late at once. CPPR gives this pessimism back; without it, hold
      shows false violations
    \item \textbf{SI}: a neighbouring net switching at the same time
      slows a transition down or speeds it up. Needs the coupling
      capacitances of \cmd{rc\_coupled}
  \end{itemize}
\end{block}
\end{frame}

\renewcommand{\assubtitle}{Running the STA}
\begin{frame}[fragile]
\frametitle{\astitle}
\begin{block}{\assubtitle}
\begin{lstlisting}[style=sh]
./scripts/sta/sta.sh -b        # batch, summary
./scripts/sta/sta.sh           # Tempus shell stays open
./scripts/sta/sta.sh -b --no-si
\end{lstlisting}
\begin{lstlisting}[style=out]
=== [tempus] 2/2 result
  check view            worst slack [ns]
  setup av_slow         ...
  setup av_typ          ...
  setup av_jtag_slow    ...
  setup av_scan_slow    ...
  hold  av_fast         ...
  hold  av_jtag_fast    ...
  hold  av_scan_fast    ...
  violators (timing + DRV): N line(s) in sta_violators.rpt
  SDF: cc_chip_top_sta_func.sdf
  SDF: cc_chip_top_sta_scan.sdf
\end{lstlisting}
  \scriptsize
  \begin{itemize}
    \item A negative slack is marked \cmd{VIOLATED}, exit code 1
  \end{itemize}
\end{block}
\end{frame}

\renewcommand{\assubtitle}{Reading the Reports}
\begin{frame}
\frametitle{\astitle}
\begin{block}{\assubtitle}
  \begin{itemize}
    \item \cmd{sta\_check\_timing.rpt} first: unconstrained endpoints,
      clocks that do not reach a flip-flop, combinational loops.
      Expected: reset pins (false path), asynchronous GPIO/UART paths
    \item \cmd{sta\_coverage.rpt}: which checks are analysed in which
      view -- a mode with almost nothing analysed means a wrong
      \cmd{set\_case\_analysis} or a renamed pin
    \item \cmd{sta\_setup/hold\_<view>.rpt}: per path the launch clock,
      the cells and wires with their delays, the capture clock and
      \emph{slack = required $-$ arrival}
    \item \cmd{sta\_violators.rpt}: also DRV --
      \cmd{max\_transition}, \cmd{max\_capacitance}, \cmd{max\_fanout}:
      slow edges are a reliability problem even if the timing is met
  \end{itemize}
\end{block}
\end{frame}

\renewcommand{\assubtitle}{Interactive Queries}
\begin{frame}[fragile]
\frametitle{\astitle}
\begin{block}{\assubtitle}
  \begin{itemize}
    \item \cmd{./scripts/sta/sta.sh} without \cmd{-b}: after the run the
      Tempus prompt stays open
  \end{itemize}
\begin{lstlisting}[style=sh]
report_timing -late  -view av_slow -to   core/cpu/...    ;# one endpoint
report_timing -early -view av_scan_fast -max_paths 5
report_timing -late  -view av_slow -path_type full_clock ;# clock path too
report_clocks
report_case_analysis                       ;# which pins are fixed
\end{lstlisting}
  \begin{itemize}
    \item \cmd{-path\_type full\_clock} shows the clock tree in the path:
      the place to look for skew in a hold violation
  \end{itemize}
\end{block}
\end{frame}

\renewcommand{\assubtitle}{What If It Fails?}
\begin{frame}
\frametitle{\astitle}
\begin{block}{\assubtitle}
  \footnotesize
  \begin{tabularx}{\linewidth}{@{}lY@{}}
    \textbf{symptom} & \textbf{fix}\\\hline
    hold, all views   & post-route hold fixing in Innovus
                        (\cmd{opt\_design -post\_route -hold}), then
                        export, extract, STA again\\
    hold, only scan   & same; check the clock gate/mux skew between
                        \cmd{tck\_i} and the core clock tree\\
    setup, func       & 800\,ns budget -- a violation means a wrong
                        constraint, not a slow path: check the clock\\
    setup, only scan  & slower TCK for SCANTEST capture, note it\\
    no paths in a view & pin names in the SDC changed in the netlist
                        (ungroup) -- \cmd{check\_timing}\\
    DRV               & \cmd{opt\_design -post\_route -drv} in Innovus\\
  \end{tabularx}
  \vspace{4pt}
  \begin{itemize}
    \item The P\&R script has no post-route optimization step yet
  \end{itemize}
\end{block}
\end{frame}

%%%%%%%%%%%%%%%%%%%%%%%%%%%%%%%%
% GLS
%%%%%%%%%%%%%%%%%%%%%%%%%%%%%%%%
\section{Post-Route GLS}
\renewcommand{\astitle}{Post-Route Gate-Level Simulation}

\renewcommand{\assubtitle}{With the Tempus SDF}
\begin{frame}[fragile]
\frametitle{\astitle}
\begin{block}{\assubtitle}
\begin{lstlisting}[style=sh]
./scripts/gatesim/gatesim.sh -b -p instr06addi   # routed netlist
./scripts/gatesim/gatesim.sh -r -p -s            # SDF max (setup)
./scripts/gatesim/gatesim.sh -r -p -s -m         # SDF min (hold)
\end{lstlisting}
  \begin{itemize}
    \item \cmd{-p}: \cmd{build\_innovus/cc\_chip\_top\_routed.v} instead
      of the Genus netlist, results in \cmd{build\_gatesim/postroute/}
    \item \cmd{-s}: \cmd{build\_sta/cc\_chip\_top\_sta\_func.sdf}, for the
      test \cmd{scan} the \cmd{\_sta\_scan.sdf}
    \item \cmd{-m}: min delays (fast corner); default max (slow corner)
    \item STA says \emph{whether} the timing is met; the GLS with SDF
      shows that the design works with these delays and that the
      constraints describe what the testbench does
  \end{itemize}
\end{block}
\end{frame}

%%%%%%%%%%%%%%%%%%%%%%%%%%%%%%%%
% Open points
%%%%%%%%%%%%%%%%%%%%%%%%%%%%%%%%
\section{Open Points}
\renewcommand{\astitle}{Open Points}

\renewcommand{\assubtitle}{VERIFY-LAB: Not Yet Run with the Real Tools}
\begin{frame}
\frametitle{\astitle}
\begin{block}{\assubtitle}
  \begin{itemize}
    \item Quantus CCL syntax and the SPEF names with three corners
      (expected \cmd{<file\_name>\_<corner>.spef})
    \item Layer names LEF vs. \cmd{qrcTechFile}: \cmd{QRC\_LAYERMAP}
      needed?
    \item \cmd{tempus -stylus} with \cmd{-files}/\cmd{-log}, the Stylus
      names \cmd{read\_netlist -top}, \cmd{delaycal\_enable\_si}
    \item \cmd{set\_analysis\_view} called twice (\cmd{xh018\_mmmc.tcl},
      then \cmd{xh018\_mmmc\_sta.tcl}): the second must win
    \item \cmd{write\_def}/\cmd{write\_netlist} options in Innovus 26.1
    \item Worst-slack summary via \cmd{report\_timing -collection}: if it
      prints \cmd{n/a}, read the reports
  \end{itemize}
\end{block}
\end{frame}

\renewcommand{\assubtitle}{X-FAB NDA}
\begin{frame}
\frametitle{\astitle}
\begin{block}{\assubtitle}
  \begin{itemize}
    \item DEF, SPEF, SDF, reports and the generated \cmd{quantus.ccl} and
      \cmd{lib.defs} contain X-FAB cell names and paths
    \item \cmd{build\_innovus/}, \cmd{build\_quantus/}, \cmd{build\_sta/}
      are in \cmd{.gitignore}: never commit them, never copy them out of
      the lab
    \item Check \cmd{git status} before every commit
  \end{itemize}
\end{block}
\end{frame}

\renewcommand{\assubtitle}{Next Steps}
\begin{frame}
\frametitle{\astitle}
\begin{block}{\assubtitle}
  \begin{itemize}
    \item Formal equivalence with Conformal LEC: RTL $\leftrightarrow$
      Genus netlist $\leftrightarrow$ Innovus netlist (scan mode
      constrained off)
    \item Physical verification: DRC, LVS, ERC with the runsets the
      X-FAB PDK provides (Pegasus, PVS, Calibre or Assura); GDS merge
      with standard cells, pads and SRAM macros
    \item Before tape-out: power ring/stripes, pad order (TECH-06), and
      this STA again on the final layout
  \end{itemize}
\end{block}
\end{frame}

%%%%%%%%%%%%%%%%%%%%%%%%%%%%%%%%
% Q&A
%%%%%%%%%%%%%%%%%%%%%%%%%%%%%%%%
\section{Q\&A}
\renewcommand{\astitle}{Q\&A}
\begin{frame}
\frametitle{\astitle}
\vfill
\centering
\begin{Huge}\textcolor{red}{Questions?}\end{Huge}
\vfill
\end{frame}

\end{document}
